\subsection{Task, researcher state and epistemic state}
A scoped inquiry is
\begin{equation}\mathcal P=(O,\mathcal Q,\Gamma_0,\mathcal E_0,\mathcal B,\Lambda),\end{equation}
where $O$ is the target object, $\mathcal Q$ the registered scientific questions or quantities of interest, $\Gamma_0$ the context/observation boundary, $\mathcal E_0$ the evidence available at the declared cutoff, $\mathcal B$ the resource budget and $\Lambda$ the blocking epistemic invariants.

The functional researcher state is
\begin{equation}\mathfrak R_t=(K_t,\mathcal Z_t,\Omega_t,\Pi_t,\mathcal G_t,\mathcal M_t,\mathcal X_t,\mathcal R_t),\end{equation}
separating canonical scientific state, explanatory models, method skills, executive policy, research agenda, metacognition, experience/transfer history and resources. The epistemic state is
\begin{equation}K_t=(\mathcal A_t,\mathcal T_t,\mathcal V_t,\mathcal E_t,\mathcal U_t,\mathcal O_t,\mathcal F_t,\mathcal H_t^{-},\mathcal S_t,\mathcal G_t^K),\end{equation}
with Knowledge Atlas charts, typed transitions, survivors, evidence/provenance, uncertainty and identified sets, obstructions/residuals, recursive fibers, immutable negative history, saturation state and protected governance identity.

% RAKL_V3_MANUSCRIPT_UPDATE_20260811

Orion places this epistemic state inside a broader persistent value state $R_t$. We write
\begin{equation}
K_t=\pi_{\mathrm{epi}}(R_t),
\end{equation}
where $\pi_{\mathrm{epi}}$ is the authority-preserving epistemic projection. The remaining v3 views contain immutable task episodes, versioned lessons, operational tools, failure diagnoses, strategy/expertise abstractions, vector-saturation state and an optional archive of Self-Orion variants. The global substrate is a typed relational object, not a claim that all of these views form one mathematical lattice. Specialized stores remain the semantic and authority owners of their objects.

For a replaceable driver with fixed weights $\theta$, one task turn can be written schematically as
\begin{equation}
(S_t,\tau_t)=\operatorname{Driver}_{\theta}(P_t,R_t),
\qquad
R_{t+1}=\operatorname{Learn}(R_t,\tau_t).
\end{equation}
The learning map changes external state rather than silently changing the evidential meaning of $K_t$. In particular, experience-derived routing priors may alter which admissible operator is attempted first, but scientific promotion still passes through the ordinary evidence and governance transitions of the epistemic projection.

\subsection{What a local scientific chart must contain}

A local chart is not simply a paper abstract placed in a database. It is the smallest context-scoped scientific object that can participate in comparison. At minimum it identifies the scientific object or system, the proposition, the quantity or relation asserted, the population/regime, the observation or measurement operator, units and transforms where relevant, uncertainty/error semantics, source evidence, assumptions, falsifiers, and the authority coordinates currently licensed. Missing fields remain missing rather than being supplied from neighboring papers.

This requirement is intentionally stricter than ordinary information extraction. A sentence such as ``X increases Y'' cannot be compared safely until ``X,'' ``Y,'' intervention status, time scale and population are resolved. In one source the statement may describe an observational association; in another it may describe an intervention; in a third it may be a fitted model derivative. The surface predicate is shared while the scientific transition type is not. Projection is the step that makes this difference explicit.

Charts can be multi-resolution. A coarse chart can be useful for search while a fine chart preserves the measurement details required for analysis. The coarse object does not replace the fine one; it carries a reconstruction pointer and a declaration of erased coordinates. This provides a formal place for summaries without making them canonical truth.

Overlap is also typed. Two charts overlap only on the coordinates relevant to the comparison. A population overlap does not imply a measurement overlap, and a shared variable name does not imply a shared operationalization. Restriction maps therefore state which coordinates are being aligned or forgotten. When such maps are approximate, their error enters the transition certificate.

The chart idea is especially useful for cross-disciplinary synthesis. Different fields often carve the same phenomenon at different scales, with different measurement conventions and different explanatory targets. A forced global ontology can erase those distinctions. Local charts allow each field's valid objects to remain native while explicit maps represent the conditions under which information can move between them. The target is controlled interoperability, not a universal vocabulary.

Finally, chart quality and source prestige are distinct. A highly cited source can project to an underspecified chart if the required context is absent. A modest source can provide a strong local measurement under a narrow regime. The authority system evaluates the projected object and evidence, not a venue label alone.

\subsection{Contextual charts, transitions and gluing}
A local scientific chart is
\begin{equation}C_i=(U_i,\phi_i,\gamma_i,e_i,\alpha_i),\end{equation}
where $\gamma_i$ can include population, scale, regime, units, observation model, assumptions, intervention, time/cutoff and QoI. Source observations are contextual projections $y_i=\pi_i(O\mid\gamma_i)$ rather than globally competing whole-object answers.

For overlapping charts,
\begin{equation}T_{ij}^{\tau,\sigma}:\phi_i(U_i\cap U_j)\rightarrow\phi_j(U_i\cap U_j)\end{equation}
is a typed, scoped transition carrying evidence, assumptions and uncertainty semantics. Composition is licensed only when interfaces, relation types, scopes, invariants, lineage and error-composition rules are compatible. Graph connectivity alone does not create endpoint scientific authority.

At comparison layer $\ell$, a contradiction requires overlap, alignment and incompatibility:
\begin{equation}\operatorname{Contradict}_{\ell}(C_i,C_j)=\operatorname{Overlap}(C_i,C_j)\land\operatorname{Align}_{\ell}(C_i,C_j)\land\neg\operatorname{Compatible}_{\ell}(C_i,C_j).\end{equation}
The gluing operator may return a global formalism, a plural atlas, or an obstructed/identified set. Pairwise compatibility is therefore not silently promoted to unique global identification.

\begin{figure}[t]
\centering
\resizebox{0.98\textwidth}{!}{%
\begin{tikzpicture}[>=Latex,every node/.style={font=\sffamily\footnotesize,align=center},chart/.style={draw,rounded corners=1.5pt,minimum width=29mm,minimum height=17mm,fill=black!3},target/.style={draw,double,rounded corners=1.5pt,minimum width=28mm,minimum height=14mm,fill=black!7},cut/.style={draw,dashed,rounded corners=1.5pt,minimum width=30mm,minimum height=13mm},trans/.style={->,line width=0.6pt},strongpath/.style={->,line width=0.95pt},weak/.style={->,dashed,line width=0.55pt}]
\node[chart] (c1) at (0,1.55) {$C_1$\\population A\\representation $\phi_1$};
\node[chart] (c2) at (4.05,1.55) {$C_2$\\population A\\representation $\phi_2$};
\node[chart] (c3) at (0,-1.45) {$C_3$\\population B\\measurement $h_3$};
\node[chart] (c4) at (4.05,-1.45) {$C_4$\\population A\\mechanistic view};
\node[target] (tau) at (8.45,0.9) {target $\tau$\\$(q,\alpha^*,\gamma)$};
\node[cut] (cut) at (8.45,-1.65) {missing prerequisite\\epistemic cut $B^*_{\tau}$};
\draw[trans] (c1) -- node[above]{typed $T_{12}^{\tau,\sigma}$} (c2);
\draw[weak] (c1) -- node[left,xshift=-1mm]{context\\translation} (c3);
\draw[trans] (c2) -- node[right,xshift=1.5mm,text width=23mm]{ancestry /\\derivation} (c4);
\draw[dashed,line width=0.55pt] (c3.east) -- (c4.west); \node[fill=white,inner sep=2pt,font=\sffamily\scriptsize] at (2.03,-1.45) {obstruction};
\draw[strongpath] (c2.east) to[bend left=7] (tau.north west); \draw[strongpath] (c4.east) to[bend right=9] (tau.south west); \draw[weak] (cut.north) -- (tau.south);
\node[font=\sffamily\scriptsize,text width=45mm,anchor=north west] at (-1.35,-2.95) {Local papers, datasets and derived models are contextual charts rather than globally competing whole objects. Pairwise compatibility is not enough for a unique global theory.};
\node[font=\sffamily\scriptsize,text width=49mm,anchor=north west] at (4.55,-2.95) {The active LLM context follows the smallest authority-valid support corridor to the target. If every admissible route crosses an unresolved prerequisite, that prerequisite becomes an epistemic cut and a new fiber.};
\end{tikzpicture}}
\caption{\textbf{Contextual Knowledge Atlas and target-conditioned inquiry.} Papers, datasets and models become local scientific charts. Typed transitions establish which comparisons and derivations are licensed; unresolved incompatibility remains an obstruction. The active research problem is to find an authority-valid support structure to target $\tau$.}
\label{fig:atlas}
\end{figure}

\subsection{Four local-to-global questions that must not be collapsed}

The atlas vocabulary is meant to separate four questions that are often compressed into the word ``consistency.'' The first is \emph{local compatibility}: after aligning the relevant population, scale, measurement operator and regime, do two charts make claims that can jointly hold on their overlap? The second is \emph{path coherence}: if several typed transitions are composed, do their scopes, assumptions and error semantics remain compatible along the path? The third is \emph{global existence}: is there at least one global object whose restrictions agree with the retained local charts? The fourth is \emph{global identifiability}: if such global objects exist, does the evidence distinguish one of them?

These predicates can fail independently. Pairwise-compatible charts can participate in a cycle whose composed transition is inconsistent. A global model can exist but remain non-unique. A unique mathematical gluing can exist while the evidence supporting one chart is too weak for the scientific authority requested by the target. Conversely, two high-authority local results can be genuinely incompatible because a hidden context coordinate prevents common gluing. Treating all four cases as ``contradiction'' loses the information needed to choose the next experiment.

The transition registry therefore records relation type, direction, domain, codomain, scope and approximation. A coordinate transform, coarse-graining map, causal intervention relation and analogy are not interchangeable arrows. If a path mixes them, the composed path must state what kind of endpoint claim it licenses. For example, an exact unit conversion followed by an approximate model reduction does not yield an exact end-to-end identity. The remainder term from the reduction remains part of the path certificate.

Obstructions are useful because they localize failure. An obstruction may indicate a missing overlap witness, an incompatible context, a failed limiting relation, an observation model that cannot transport, or a model family that lacks a required degree of freedom. Recent sheaf-theoretic work on scientific theory shift similarly uses transport failure and obstruction as signals that a representational language may need extension \cite{olivieri2026sheaf}. \RAKL{} incorporates that logic while also asking what evidence authority the proposed extension would require.

A practical consequence is that ``merge'' is never the default response to agreement. Two charts can remain separate even when compatible if their contexts matter for downstream transport or if collapsing them would destroy provenance. Gluing is a scientific operation with an information-loss profile. When a global summary is constructed, the atlas should preserve the local charts and the gluing witness so later evidence can reopen only the affected region instead of invalidating an undifferentiated global narrative.

\subsection{Multi-axis scientific authority}
For claim $c$, scientific authority is represented as
\begin{equation}\alpha(c)=(G_c,R_c,M_c,I_c,D_c),\end{equation}
where the coordinates are scoped certificate sets for grounding/provenance, representation/relation, mechanism ancestry, identification/bounding and decision/QoI usability. Under compatible scope, the induced order is coordinatewise. It is a partial order, not automatically a mathematical lattice: incompatible assumptions, populations or observation operators can obstruct a join.

Cross-axis non-escalation is explicit:
\begin{equation}R\not\Rightarrow M,\qquad D\not\Rightarrow M,\qquad M\not\Rightarrow I.\end{equation}
Thus predictive or observational success does not by itself identify a mechanism; mechanism plausibility does not establish point identification; provenance and citation multiplicity do not create truth or independent evidence.
